%! Conic Sections — Unit 4, Lesson 4.6 — Warm-Up (Answer Key)
\documentclass[10pt]{article}
\usepackage{precalculus-article}
\usepackage{precalculus-key}   % pulls in -boxes; do NOT also load -boxes
\newcommand{\TallMath}[1]{$\displaystyle #1\rule[-1.4em]{0pt}{3.2em}$}

\begin{document}

\pageheader{Unit 4, Lesson 4.6}{Warm-Up --- Answer Key}
\namedateperiod

\vspace{0.12in}

\begin{notesbox}{Spiral Review}

\textbf{1.}\quad Complete the square to rewrite $x^2 - 6x$ in the form
$(x - h)^2 + c$. What is $h$, and what constant $c$ must you add?

\ansline{$(x-3)^2 - 9$; \quad $h = 3$; \quad $c = -9$ (i.e., subtract 9 to balance)}

\vspace{0.06in}

\textbf{2.}\quad A circle is centered at $(2, -3)$ with radius $5$.
Write its equation in standard form.

\ansline{$(x-2)^2 + (y+3)^2 = 25$}

\vspace{0.06in}

\textbf{3.}\quad Consider the equation
$\dfrac{(x-1)^2}{9} + \dfrac{(y+2)^2}{4} = 1$.

\vspace{4pt}
(a)\quad What is the center of this ellipse?
\ansline{$(1,\,-2)$}

(b)\quad What is the horizontal radius $a$?
\ansline{$a = 3$ \quad (since $a^2 = 9$)}

(c)\quad What is the vertical radius $b$?
\ansline{$b = 2$ \quad (since $b^2 = 4$)}

\end{notesbox}

\begin{teachernote}
Common errors to watch: students may report $a = 9$ (forgetting to take the
square root) or read the center as $(−1, 2)$ (flipping the signs of $h$ and
$k$). Address both before moving into the lesson proper.
\end{teachernote}

\end{document}
